\section{Rings of fractions}
\label{section:0.1}

\setcounter{subsection}{-1}
\subsection{Rings and algebras}
\label{subsection:0.1.0}

\oldpage[0\textsubscript{I}]{11}
\begin{env}[1.0.1]
\label{0.1.0.1}
Every ring considered in this Treatise has a \emph{unit element};
every module over such a ring is assumed to be \emph{unital}; ring homomorphisms will always be assumed to \emph{map the unit element to the unit element}; unless expressly stated otherwise, a subring of a ring $A$ will be assumed to \emph{contain the unit element of $A$}.
We shall consider mainly \emph{commutative} rings, and when we speak of a ring without further qualification, it will be understood that a commutative ring is meant.
If $A$ is a not necessarily commutative ring, by an $A$-module we shall always mean a \emph{left} module, unless expressly stated otherwise.
\end{env}

\begin{env}[1.0.2]
\label{0.1.0.2}
Let $A$, $B$ be two not necessarily commutative rings, and $\varphi:A\to B$ a homomorphism.
Every left (resp. right) $B$-module $M$ can be endowed with a left (resp. right) $A$-module structure by setting $a.m=\varphi(a).m$ (resp. $m.a=m.\varphi(a)$); when it is necessary to distinguish between the $A$-module and $B$-module structures on $M$, we shall denote by $M_{[\varphi]}$ the left (resp. right) $A$-module thus defined.
If $L$ is an $A$-module, a homomorphism $u:L\to M_{[\varphi]}$ is therefore a homomorphism of commutative groups such that $u(a.x)=\varphi(a).u(x)$ for $a\in A$, $x\in L$; we shall also say that it is a \emph{$\varphi$-homomorphism} $L\to M$, and that the pair $(\varphi,u)$ (or, by abuse of language, $u$) is a \emph{di-homomorphism} from $(A,L)$ to $(B,M)$.
The pairs $(A,L)$ consisting of a ring $A$ and an $A$-module $L$ therefore form a \emph{category} whose morphisms are the di-homomorphisms.
\end{env}

\begin{env}[1.0.3]
\label{0.1.0.3}
Under the hypotheses of (1.0.2), if $\mathfrak{J}$ is a left (resp. right) ideal of $A$, we shall denote by $B\mathfrak{J}$ (resp. $\mathfrak{J}B$) the left (resp. right) ideal $B\varphi(\mathfrak{J})$ (resp. $\varphi(\mathfrak{J})B$) of $B$ generated by $\varphi(\mathfrak{J})$; it is also the image of the canonical homomorphism $B\otimes_A\mathfrak{J}\to B$ (resp. $\mathfrak{J}\otimes_A B\to B$) of left (resp. right) $B$-modules.
\end{env}

\begin{env}[1.0.4]
\label{0.1.0.4}
If $A$ is a (commutative) ring and $B$ a not necessarily commutative ring, giving an \emph{$A$-algebra} structure on $B$ is equivalent to giving a ring homomorphism $\varphi:A\to B$ such that $\varphi(A)$ is contained in the centre of $B$.
For every ideal $\mathfrak{J}$ of $A$, $\mathfrak{J}B=B\mathfrak{J}$ is then a two-sided ideal of $B$, and for every $B$-module $M$, $\mathfrak{J}M$ is then a $B$-module equal to $(B\mathfrak{J})M$.
\end{env}

\begin{env}[1.0.5]
\label{0.1.0.5}
We shall not go back over the notions of \emph{module of finite type} and of (commutative) \emph{algebra of finite type}; to say that an $A$-module $M$ is of finite type means that there exists

\oldpage[0\textsubscript{I}]{12}
an exact sequence $A^p\to M\to 0$.
An $A$-module $M$ is said to admit a \emph{finite presentation} if it is isomorphic to the cokernel of a homomorphism $A^p\to A^q$, in other words if there exists an exact sequence $A^p\to A^q\to M\to 0$.
Note that over a \emph{Noetherian} ring $A$, every $A$-module of finite type admits a finite presentation.

Recall that an $A$-algebra $B$ is said to be \emph{integral over $A$} if every element of $B$ is a root in $B$ of a monic polynomial with coefficients in $A$; it amounts to the same to say that every element of $B$ is contained in a subalgebra of $B$ which is an \emph{$A$-module of finite type}.
When this is so, and $B$ is commutative, the subalgebra of $B$ generated by a finite subset of $B$ is an $A$-module of finite type; for the commutative algebra $B$ to be integral and of finite type over $A$, it is therefore necessary and sufficient that $B$ be an $A$-module of finite type; one then also says that $B$ is a \emph{finite integral} $A$-algebra (or simply a \emph{finite} one, if no confusion results).
Observe that in these definitions the homomorphism $A\to B$ defining the $A$-algebra structure is not assumed to be injective.
\end{env}

\begin{env}[1.0.6]
\label{0.1.0.6}
An \emph{integral} ring is a ring in which the product of a finite family of elements $\neq 0$ is $\neq 0$; it amounts to the same to say that in such a ring one has $0\neq 1$ and the product of two elements $\neq 0$ is nonzero.
A \emph{prime} ideal of a ring $A$ is an ideal $\mathfrak{p}$ such that $A/\mathfrak{p}$ is integral; this therefore implies $\mathfrak{p}\neq A$.
For a ring $A$ to have at least one prime ideal, it is necessary and sufficient that $A\neq\{0\}$.
\end{env}

\begin{env}[1.0.7]
\label{0.1.0.7}
A \emph{local} ring is a ring $A$ in which there exists only one maximal ideal, which is then the complement of the set of invertible elements and contains all the ideals $\neq A$.
If $A$ and $B$ are two local rings, and $\mathfrak{m}$ and $\mathfrak{n}$ their respective maximal ideals, a homomorphism $\varphi:A\to B$ is said to be \emph{local} if $\varphi(\mathfrak{m})\subset\mathfrak{n}$ (or, what amounts to the same, if $\varphi^{-1}(\mathfrak{n})=\mathfrak{m}$).
By passage to the quotients, such a homomorphism then defines a monomorphism of the residue field $A/\mathfrak{m}$ into the residue field $B/\mathfrak{n}$.
The composite of two local homomorphisms is a local homomorphism.
\end{env}

\subsection{Radical of an ideal. Nilradical and radical of a ring}
\label{subsection:0.1.1}

\begin{env}[1.1.1]
\label{0.1.1.1}
Let $\mathfrak{a}$ be an ideal of a ring $A$; the \emph{radical} of $\mathfrak{a}$, denoted $\mathfrak{r}(\mathfrak{a})$, is the set of $x\in A$ such that $x^n\in\mathfrak{a}$ for at least one integer $n>0$; it is an ideal containing $\mathfrak{a}$.
We have $\mathfrak{r}(\mathfrak{r}(\mathfrak{a}))=\mathfrak{r}(\mathfrak{a})$; the relation $\mathfrak{a}\subset\mathfrak{b}$ implies $\mathfrak{r}(\mathfrak{a})\subset\mathfrak{r}(\mathfrak{b})$; the radical of a finite intersection of ideals is the intersection of their radicals.
If $\varphi$ is a homomorphism from a ring $A'$ to $A$, we have $\mathfrak{r}(\varphi^{-1}(\mathfrak{a}))=\varphi^{-1}(\mathfrak{r}(\mathfrak{a}))$ for every ideal $\mathfrak{a}\subset A$.
An ideal is the radical of an ideal if and only if it is an intersection of prime ideals.
The radical of an ideal $\mathfrak{a}$ is the intersection of the prime ideals which are \emph{minimal} among those containing $\mathfrak{a}$; if $A$ is Noetherian, these minimal prime ideals are finite in number.

The radical of the ideal $(0)$ is also called the \emph{nilradical} of $A$; it is the set $\mathfrak{N}$ of nilpotent elements of $A$.
The ring $A$ is said to be \emph{reduced} if $\mathfrak{N}=(0)$; for every ring $A$, the quotient $A/\mathfrak{N}$ of $A$ by its nilradical is a reduced ring.
\end{env}

\begin{env}[1.1.2]
\label{0.1.1.2}
Recall that the \emph{radical} $\mathfrak{R}(A)$ of a (not necessarily commutative) ring $A$ is the intersection of the maximal left ideals of $A$ (and also the intersection of the maximal right ideals).
The radical of $A/\mathfrak{R}(A)$ is $(0)$.
\end{env}

\subsection{Modules and rings of fractions}
\label{subsection:0.1.2}

\oldpage[0\textsubscript{I}]{13}
\begin{env}[1.2.1]
\label{0.1.2.1}
A subset $S$ of a ring $A$ is said to be \emph{multiplicative} if $1\in S$ and the product of two elements of $S$ is in $S$.
The examples which will be the most important in what follows are: $1^\circ$ the set $S_f$ of the powers $f^n$ ($n\geq 0$) of an element $f\in A$; $2^\circ$ the complement $A-\mathfrak{p}$ of a \emph{prime} ideal $\mathfrak{p}$ of $A$.
\end{env}

\begin{env}[1.2.2]
\label{0.1.2.2}
Let $S$ be a multiplicative subset of a ring $A$, and $M$ an $A$-module; in the set $M\times S$, the relation between pairs $(m_1,s_1)$, $(m_2,s_2)$:
\[
  \text{``there exists $s\in S$ such that $s(s_1m_2-s_2m_1)=0$''}
\]
is an equivalence relation.
We denote by $S^{-1}M$ the quotient set of $M\times S$ by this relation, and by $m/s$ the canonical image in $S^{-1}M$ of the pair $(m,s)$; the map $i_M^S:m\mapsto m/1$ (also denoted $i^S$) is called the \emph{canonical map} from $M$ to $S^{-1}M$.
This map is in general neither injective nor surjective; its kernel is the set of $m\in M$ for which there exists an $s\in S$ such that $sm=0$.

On $S^{-1}M$ one defines an additive group law by taking
\[
  (m_1/s_1)+(m_2/s_2)=(s_2m_1+s_1m_2)/(s_1s_2)
\]
(one checks that this is indeed independent of the expressions of the elements of $S^{-1}M$ considered).
On $S^{-1}A$ one further defines a multiplicative law by taking $(a_1/s_1)(a_2/s_2)=(a_1a_2)/(s_1s_2)$, and finally an external law on $S^{-1}M$, having $S^{-1}A$ as set of operators, by setting $(a/s)(m/s')=(am)/(ss')$.
One checks in this way that $S^{-1}A$ is endowed with a ring structure (called the \emph{ring of fractions of $A$ with denominators in $S$}) and $S^{-1}M$ with an $S^{-1}A$-module structure (called the \emph{module of fractions of $M$ with denominators in $S$}); for every $s\in S$, $s/1$ is invertible in $S^{-1}A$, its inverse being $1/s$.
The canonical map $i_A^S$ (resp. $i_M^S$) is a ring homomorphism (resp. a homomorphism of $A$-modules, $S^{-1}M$ being regarded as an $A$-module by means of the homomorphism $i_A^S:A\to S^{-1}A$).
\end{env}

\begin{env}[1.2.3]
\label{0.1.2.3}
If $S_f=\{f^n\}_{n\geq 0}$ for some $f\in A$, we write $A_f$ and $M_f$ instead of $S_f^{-1}A$ and $S_f^{-1}M$; when $A_f$ is regarded as an algebra over $A$, we may write $A_f=A[1/f]$.
$A_f$ is isomorphic to the quotient algebra $A[T]/(fT-1)A[T]$.
When $f=1$, $A_f$ and $M_f$ are canonically identified with $A$ and $M$; if $f$ is nilpotent, $A_f$ and $M_f$ are reduced to $0$.

When $S=A-\mathfrak{p}$, where $\mathfrak{p}$ is a prime ideal of $A$, we write $A_{\mathfrak{p}}$ and $M_{\mathfrak{p}}$ instead of $S^{-1}A$ and $S^{-1}M$; $A_{\mathfrak{p}}$ is a \emph{local ring} whose maximal ideal $\mathfrak{q}$ is generated by $i_A^S(\mathfrak{p})$, and we have $(i_A^S)^{-1}(\mathfrak{q})=\mathfrak{p}$; by passage to the quotients, $i_A^S$ gives a monomorphism of the integral ring $A/\mathfrak{p}$ into the field $A_{\mathfrak{p}}/\mathfrak{q}$, which is identified with the field of fractions of $A/\mathfrak{p}$.
\end{env}

\begin{env}[1.2.4]
\label{0.1.2.4}
The ring of fractions $S^{-1}A$ and the canonical homomorphism $i_A^S$ are a solution of a \emph{universal mapping problem}: every homomorphism $u$ from $A$ to a ring $B$ such that $u(S)$ consists of elements \emph{invertible} in $B$ factorizes in a unique way as
\[
  u:A\xrightarrow{i_A^S}S^{-1}A\xrightarrow{u^*}B
\]

\oldpage[0\textsubscript{I}]{14}
where $u^*$ is a ring homomorphism.
Under the same hypotheses, let $M$ be an $A$-module, $N$ a $B$-module, and $v:M\to N$ a homomorphism of $A$-modules (for the $B$-module structure on $N$ defined by $u:A\to B$); then $v$ factorizes in a unique way as
\[
  v:M\xrightarrow{i_M^S}S^{-1}M\xrightarrow{v^*}N
\]
where $v^*$ is a homomorphism of $S^{-1}A$-modules (for the $S^{-1}A$-module structure on $N$ defined by $u^*$).
\end{env}

\begin{env}[1.2.5]
\label{0.1.2.5}
One defines a canonical isomorphism $S^{-1}A\otimes_A M\xrightarrow{\sim}S^{-1}M$ of $S^{-1}A$-modules by mapping the element $(a/s)\otimes m$ to the element $(am)/s$, the inverse isomorphism mapping $m/s$ to $(1/s)\otimes m$.
\end{env}

\begin{env}[1.2.6]
\label{0.1.2.6}
For every ideal $\mathfrak{a}'$ of $S^{-1}A$, $\mathfrak{a}=(i_A^S)^{-1}(\mathfrak{a}')$ is an ideal of $A$, and $\mathfrak{a}'$ is the ideal of $S^{-1}A$ generated by $i_A^S(\mathfrak{a})$, which is identified with $S^{-1}\mathfrak{a}$ (1.3.2).
The map $\mathfrak{p}'\mapsto(i_A^S)^{-1}(\mathfrak{p}')$ is an isomorphism, for the order structures, of the set of \emph{prime} ideals of $S^{-1}A$ onto the set of prime ideals $\mathfrak{p}$ of $A$ such that $\mathfrak{p}\cap S=\varnothing$.
Moreover, the local rings $A_{\mathfrak{p}}$ and $(S^{-1}A)_{S^{-1}\mathfrak{p}}$ are then canonically (1.5.1) isomorphic.
\end{env}

\begin{env}[1.2.7]
\label{0.1.2.7}
When $A$ is an \emph{integral} ring, whose field of fractions we denote by $K$, the canonical map $i_A^S:A\to S^{-1}A$ is injective for every multiplicative subset $S$ not containing $0$, and $S^{-1}A$ is then canonically identified with a subring of $K$ containing $A$.
In particular, for every prime ideal $\mathfrak{p}$ of $A$, $A_{\mathfrak{p}}$ is a local ring containing $A$, with maximal ideal $\mathfrak{p}A_{\mathfrak{p}}$, and we have $\mathfrak{p}A_{\mathfrak{p}}\cap A=\mathfrak{p}$.
\end{env}

\begin{env}[1.2.8]
\label{0.1.2.8}
If $A$ is a \emph{reduced} ring (1.1.1), so is $S^{-1}A$: indeed, if $(x/s)^n=0$ for $x\in A$, $s\in S$, this means that there exists $s'\in S$ such that $s'x^n=0$, whence $(s'x)^n=0$, which by hypothesis implies $s'x=0$, hence $x/s=0$.
\end{env}

\subsection{Functorial properties}
\label{subsection:0.1.3}

\begin{env}[1.3.1]
\label{0.1.3.1}
Let $M$, $N$ be two $A$-modules, and $u$ an $A$-homomorphism $M\to N$.
If $S$ is a multiplicative subset of $A$, one defines an $S^{-1}A$-homomorphism $S^{-1}M\to S^{-1}N$, denoted $S^{-1}u$, by setting $(S^{-1}u)(m/s)=u(m)/s$; if $S^{-1}M$ and $S^{-1}N$ are canonically identified with $S^{-1}A\otimes_A M$ and $S^{-1}A\otimes_A N$ (1.2.5), $S^{-1}u$ is identified with $1\otimes u$.
If $P$ is a third $A$-module and $v$ an $A$-homomorphism $N\to P$, we have $S^{-1}(v\circ u)=(S^{-1}v)\circ(S^{-1}u)$; in other words, $S^{-1}M$ is a \emph{covariant functor in $M$}, from the category of $A$-modules to that of $S^{-1}A$-modules ($A$ and $S$ being fixed).
\end{env}

\begin{env}[1.3.2]
\label{0.1.3.2}
The functor $S^{-1}M$ is \emph{exact}; in other words, if the sequence
\[
  M\xrightarrow{u}N\xrightarrow{v}P
\]
is exact, so is the sequence
\[
  S^{-1}M\xrightarrow{S^{-1}u}S^{-1}N\xrightarrow{S^{-1}v}S^{-1}P.
\]
In particular, if $u:M\to N$ is injective (resp. surjective), so is $S^{-1}u$;

\oldpage[0\textsubscript{I}]{15}
if $N$ and $P$ are two submodules of $M$, $S^{-1}N$ and $S^{-1}P$ are canonically identified with submodules of $S^{-1}M$, and we have
\[
  S^{-1}(N+P)=S^{-1}N+S^{-1}P
  \quad\text{and}\quad
  S^{-1}(N\cap P)=(S^{-1}N)\cap(S^{-1}P).
\]
\end{env}

\begin{env}[1.3.3]
\label{0.1.3.3}
Let $(M_\alpha,\varphi_{\beta\alpha})$ be an inductive system of $A$-modules; then $(S^{-1}M_\alpha,S^{-1}\varphi_{\beta\alpha})$ is an inductive system of $S^{-1}A$-modules.
Expressing the $S^{-1}M_\alpha$ and $S^{-1}\varphi_{\beta\alpha}$ as tensor products (1.2.5 and 1.3.1), it follows from the fact that the operations of tensor product and of inductive limit commute that there is a canonical isomorphism
\[
  S^{-1}\varinjlim M_\alpha\xrightarrow{\sim}\varinjlim S^{-1}M_\alpha,
\]
which one also expresses by saying that the functor $S^{-1}M$ (in $M$) \emph{commutes with inductive limits}.
\end{env}

\begin{env}[1.3.4]
\label{0.1.3.4}
Let $M$, $N$ be two $A$-modules; there exists a canonical \emph{functorial} isomorphism (in $M$ and $N$)
\[
  (S^{-1}M)\otimes_{S^{-1}A}(S^{-1}N)
  \xrightarrow{\sim}
  S^{-1}(M\otimes_A N)
\]
which maps $(m/s)\otimes(n/t)$ to $(m\otimes n)/st$.
\end{env}

\begin{env}[1.3.5]
\label{0.1.3.5}
Similarly, there is a \emph{functorial} homomorphism (in $M$ and $N$)
\[
  S^{-1}\operatorname{Hom}_A(M,N)
  \longrightarrow
  \operatorname{Hom}_{S^{-1}A}(S^{-1}M,S^{-1}N)
\]
which maps $u/s$ to the homomorphism $m/t\mapsto u(m)/st$.
When $M$ has a \emph{finite presentation}, the preceding homomorphism is an \emph{isomorphism}: this is immediate when $M$ is of the form $A^r$, and one passes from there to the general case by starting from the exact sequence $A^p\to A^q\to M\to 0$ and using the exactness of the functor $S^{-1}M$ and the left exactness of the functor $\operatorname{Hom}_A(M,N)$ in $M$.
Note that this case always occurs when $A$ is \emph{Noetherian} and the $A$-module $M$ is \emph{of finite type}.
\end{env}

\subsection{Change of multiplicative subset}
\label{subsection:0.1.4}

\begin{env}[1.4.1]
\label{0.1.4.1}
Let $S$, $T$ be two multiplicative subsets of a ring $A$ such that $S\subset T$; there exists a canonical homomorphism $\rho_A^{T,S}$ (or simply $\rho^{T,S}$) from $S^{-1}A$ to $T^{-1}A$, which maps the element denoted $a/s$ of $S^{-1}A$ to the element denoted $a/s$ in $T^{-1}A$; we have
\[
  i_A^T=\rho_A^{T,S}\circ i_A^S.
\]
For every $A$-module $M$, there likewise exists an $S^{-1}A$-linear map from $S^{-1}M$ to $T^{-1}M$ (the latter being regarded as an $S^{-1}A$-module by means of the homomorphism $\rho_A^{T,S}$), which maps the element $m/s$ of $S^{-1}M$ to the element $m/s$ of $T^{-1}M$; this map is denoted $\rho_M^{T,S}$, or simply $\rho^{T,S}$, and we again have
\[
  i_M^T=\rho_M^{T,S}\circ i_M^S;
\]
under the canonical identification (1.2.5), $\rho_M^{T,S}$ is identified with $\rho_A^{T,S}\otimes 1$.
The homomorphism $\rho_M^{T,S}$ is a \emph{functorial morphism} (or natural transformation) from the functor $S^{-1}M$ to the functor $T^{-1}M$; in other words, the diagram
\[
\begin{tikzcd}[column sep=large,row sep=large]
S^{-1}M \arrow[r,"S^{-1}u"] \arrow[d,"\rho_M^{T,S}"'] & S^{-1}N \arrow[d,"\rho_N^{T,S}"] \\
T^{-1}M \arrow[r,"T^{-1}u"] & T^{-1}N
\end{tikzcd}
\]

\oldpage[0\textsubscript{I}]{16}
is commutative for every homomorphism $u:M\to N$; note moreover that $T^{-1}u$ is entirely determined by $S^{-1}u$, since for $m\in M$ and $t\in T$ we have
\[
  (T^{-1}u)(m/t)=(t/1)^{-1}\rho^{T,S}\bigl((S^{-1}u)(m/1)\bigr).
\]
\end{env}

\begin{env}[1.4.2]
\label{0.1.4.2}
With the same notation, for two $A$-modules $M$, $N$, the diagrams (cf.\ (1.3.4) and (1.3.5))
\[
\begin{tikzcd}[column sep=large,row sep=large]
(S^{-1}M)\otimes_{S^{-1}A}(S^{-1}N) \arrow[r,"\sim"] \arrow[d]
  & S^{-1}(M\otimes_A N) \arrow[d]
  \\
(T^{-1}M)\otimes_{T^{-1}A}(T^{-1}N) \arrow[r,"\sim"]
  & T^{-1}(M\otimes_A N)
\end{tikzcd}
\]
\[
\begin{tikzcd}[column sep=large,row sep=large]
S^{-1}\operatorname{Hom}_A(M,N) \arrow[r] \arrow[d]
  & \operatorname{Hom}_{S^{-1}A}(S^{-1}M,S^{-1}N) \arrow[d] \\
T^{-1}\operatorname{Hom}_A(M,N) \arrow[r]
  & \operatorname{Hom}_{T^{-1}A}(T^{-1}M,T^{-1}N)
\end{tikzcd}
\]
are commutative.
\end{env}

\begin{env}[1.4.3]
\label{0.1.4.3}
There is an important case in which the homomorphism $\rho^{T,S}$ is \emph{bijective}, namely when every element of $T$ is a divisor of an element of $S$; one then identifies the modules $S^{-1}M$ and $T^{-1}M$ by means of $\rho^{T,S}$.
We say that $S$ is \emph{saturated} if every divisor in $A$ of an element of $S$ is in $S$; by replacing $S$ by the set $T$ of all divisors of the elements of $S$ (a set which is multiplicative and saturated), one sees that one can always, if one wishes, restrict oneself to considering modules of fractions $S^{-1}M$ where $S$ is saturated.
\end{env}

\begin{env}[1.4.4]
\label{0.1.4.4}
If $S$, $T$, $U$ are three multiplicative subsets of $A$ such that $S\subset T\subset U$, we have
\[
  \rho^{U,S}=\rho^{U,T}\circ\rho^{T,S}.
\]
\end{env}

\begin{env}[1.4.5]
\label{0.1.4.5}
Consider an \emph{increasing filtered family} $(S_\alpha)$ of multiplicative subsets of $A$ (we write $\alpha\leq\beta$ for $S_\alpha\subset S_\beta$), and let $S$ be the multiplicative subset $\bigcup_\alpha S_\alpha$; set
\[
  \rho_{\beta\alpha}=\rho_A^{S_\beta,S_\alpha}\quad\text{for }\alpha\leq\beta;
\]
by (1.4.4), the homomorphisms $\rho_{\beta\alpha}$ define a ring $A'$, the \emph{inductive limit} of the inductive system of rings $(S_\alpha^{-1}A,\rho_{\beta\alpha})$.
Let $\rho_\alpha$ be the canonical map $S_\alpha^{-1}A\to A'$, and set $\varphi_\alpha=\rho_A^{S,S_\alpha}$; since $\varphi_\alpha=\varphi_\beta\circ\rho_{\beta\alpha}$ for $\alpha\leq\beta$ by (1.4.4), one can define in a unique way a homomorphism $\varphi:A'\to S^{-1}A$ such that the diagram
\[
\begin{tikzcd}[column sep=huge,row sep=large]
  & S_\alpha^{-1}A
      \arrow[ddl,"\rho_\alpha"']
      \arrow[d,"\rho_{\beta\alpha}"]
      \arrow[ddr,"\varphi_\alpha"] & \\
  & S_\beta^{-1}A
      \arrow[dl,"\rho_\beta"]
      \arrow[dr,"\varphi_\beta"'] & \\
A' \arrow[rr,"\varphi"'] & & S^{-1}A
\end{tikzcd}
\qquad(\alpha\leq\beta)
\]
is commutative.
In fact, $\varphi$ is an \emph{isomorphism}: indeed, it is immediate from the construction that $\varphi$ is surjective.
On the other hand, if $\rho_\alpha(a/s_\alpha)\in A'$ is such that $\varphi(\rho_\alpha(a/s_\alpha))=0$, this means that $a/s_\alpha=0$ in $S^{-1}A$, that is, that there exists $s\in S$ such that $sa=0$; but there is a $\beta\geq\alpha$ such that $s\in S_\beta$, and consequently, since $\rho_\alpha(a/s_\alpha)=\rho_\beta(sa/ss_\alpha)=0$, we see that $\varphi$ is injective.
The case of an $A$-module $M$ is treated in the same way, and we have thus defined canonical isomorphisms
\[
  \varinjlim S_\alpha^{-1}A\xrightarrow{\sim}
  (\varinjlim S_\alpha)^{-1}A,
  \qquad
  \varinjlim S_\alpha^{-1}M\xrightarrow{\sim}
  (\varinjlim S_\alpha)^{-1}M,
\]
the second being \emph{functorial in $M$}.
\end{env}

\oldpage[0\textsubscript{I}]{17}
\begin{env}[1.4.6]
\label{0.1.4.6}
Let $S_1$, $S_2$ be two multiplicative subsets of $A$; then $S_1S_2$ is also a multiplicative subset of $A$.
Denote by $S_2'$ the canonical image of $S_2$ in the ring $S_1^{-1}A$, which is a multiplicative subset of that ring.
For every $A$-module $M$, there then exists a functorial isomorphism
\[
  S_2'^{-1}(S_1^{-1}M)\xrightarrow{\sim}(S_1S_2)^{-1}M
\]
which maps $(m/s_1)/(s_2/1)$ to the element $m/(s_1s_2)$.
\end{env}

\subsection{Change of ring}
\label{subsection:0.1.5}

\begin{env}[1.5.1]
\label{0.1.5.1}
Let $A$, $A'$ be two rings, $\varphi$ a homomorphism $A'\to A$, and $S$ (resp. $S'$) a multiplicative subset of $A$ (resp. $A'$) such that $\varphi(S')\subset S$; the composite homomorphism
\[
  A'\xrightarrow{\varphi}A\longrightarrow S^{-1}A
\]
factorizes as
\[
  A'\longrightarrow S'^{-1}A'\xrightarrow{\varphi^{S'}}S^{-1}A
\]
by (1.2.4); we have $\varphi^{S'}(a'/s')=\varphi(a')/\varphi(s')$.
If $A=\varphi(A')$ and $S=\varphi(S')$, $\varphi^{S'}$ is \emph{surjective}.
If $A'=A$ and $\varphi$ is the identity, $\varphi^{S'}$ is none other than the homomorphism $\rho_A^{S,S'}$ defined in (1.4.1).
\end{env}

\begin{env}[1.5.2]
\label{0.1.5.2}
Under the hypotheses of (1.5.1), let $M$ be an $A$-module.
There exists a canonical functorial homomorphism
\[
  \sigma:S'^{-1}(M_{[\varphi]})\longrightarrow
  (S^{-1}M)_{[\varphi^{S'}]}
\]
of $S'^{-1}A'$-modules, which maps every element $m/s'$ of $S'^{-1}(M_{[\varphi]})$ to the element $m/\varphi(s')$ of $(S^{-1}M)_{[\varphi^{S'}]}$; indeed, one checks immediately that this definition does not depend on the expression $m/s'$ of the element considered.
When $S=\varphi(S')$, the homomorphism $\sigma$ is \emph{bijective}.
When $A'=A$ and $\varphi$ is the identity, $\sigma$ is none other than the homomorphism $\rho_M^{S,S'}$ defined in (1.4.1).

When in particular one takes $M=A$, the homomorphism $\varphi$ defines an $A'$-algebra structure on $A$; $S'^{-1}(A_{[\varphi]})$ is then endowed with a ring structure, for which it is identified with $(\varphi(S'))^{-1}A$, and the homomorphism
\[
  \sigma:S'^{-1}(A_{[\varphi]})\longrightarrow S^{-1}A
\]
is a homomorphism of $S'^{-1}A'$-algebras.
\end{env}

\begin{env}[1.5.3]
\label{0.1.5.3}
Let $M$ and $N$ be two $A$-modules; by composing the homomorphisms defined in (1.3.4) and (1.5.2), one obtains a homomorphism
\[
  (S^{-1}M\otimes_{S^{-1}A}S^{-1}N)_{[\varphi^{S'}]}
  \longleftarrow
  S'^{-1}((M\otimes_A N)_{[\varphi]})
\]
which is an isomorphism when $\varphi(S')=S$.
Similarly, by composing the homomorphisms (1.3.5) and (1.5.2), one obtains a homomorphism
\[
  S'^{-1}((\operatorname{Hom}_A(M,N))_{[\varphi]})
  \longrightarrow
  (\operatorname{Hom}_{S^{-1}A}(S^{-1}M,S^{-1}N))_{[\varphi^{S'}]}
\]
which is an isomorphism when $\varphi(S')=S$ and $M$ admits a finite presentation.
\end{env}

\begin{env}[1.5.4]
\label{0.1.5.4}
Now consider an $A'$-module $N'$, and form the tensor product $N'\otimes_{A'}A_{[\varphi]}$, which can be regarded as an $A$-module by setting
\[
  a.(n'\otimes b)=n'\otimes(ab).
\]
There exists a functorial isomorphism of $S^{-1}A$-modules
\[
  \tau:(S'^{-1}N')\otimes_{S'^{-1}A'}(S^{-1}A)_{[\varphi^{S'}]}
  \xrightarrow{\sim}
  S^{-1}(N'\otimes_{A'}A_{[\varphi]})
\]
\oldpage[0\textsubscript{I}]{18}
which maps the element $(n'/s')\otimes(a/s)$ to the element
$(n'\otimes a)/(\varphi(s')s)$; indeed, one checks separately that when
$n'/s'$ (resp. $a/s$) is replaced by another expression of the same element,
$(n'\otimes a)/(\varphi(s')s)$ does not change; on the other hand, one can define an
inverse homomorphism of $\tau$ by mapping $(n'\otimes a)/s$ to the element
$(n'/1)\otimes(a/s)$: one uses the fact that
$S^{-1}(N'\otimes_{A'}A_{[\varphi]})$ is canonically isomorphic to
$(N'\otimes_{A'}A_{[\varphi]})\otimes_A S^{-1}A$ (1.2.5), hence also to
$N'\otimes_{A'}(S^{-1}A)_{[\psi]}$, where $\psi$ denotes the composite
homomorphism $a'\mapsto\varphi(a')/1$ from $A'$ to $S^{-1}A$.
\end{env}

\begin{env}[1.5.5]
\label{0.1.5.5}
If $M'$ and $N'$ are two $A'$-modules, by composing the isomorphisms (1.3.4)
and (1.5.4) one obtains an isomorphism
\[
  S'^{-1}M'\otimes_{S'^{-1}A'}S'^{-1}N'
  \otimes_{S'^{-1}A'}S^{-1}A
  \xrightarrow{\sim}
  S^{-1}(M'\otimes_{A'}N'\otimes_{A'}A).
\]
Similarly, if $M'$ admits a finite presentation, one has, by (1.3.5) and
(1.5.4), an isomorphism
\[
  \operatorname{Hom}_{S'^{-1}A'}(S'^{-1}M',S'^{-1}N')
  \otimes_{S'^{-1}A'}S^{-1}A
  \xrightarrow{\sim}
  S^{-1}(\operatorname{Hom}_{A'}(M',N')\otimes_{A'}A).
\]
\end{env}

\begin{env}[1.5.6]
\label{0.1.5.6}
Under the hypotheses of (1.5.1), let $T$ (resp. $T'$) be a second multiplicative
subset of $A$ (resp. $A'$) such that $S\subset T$ (resp. $S'\subset T'$)
and $\varphi(T')\subset T$. Then the diagram
\[
\begin{tikzcd}[column sep=huge,row sep=large]
S'^{-1}A' \arrow[r,"\varphi^{S'}"] \arrow[d,"\rho^{T',S'}"']
  & S^{-1}A \arrow[d,"\rho^{T,S}"] \\
T'^{-1}A' \arrow[r,"\varphi^{T'}"']
  & T^{-1}A
\end{tikzcd}
\]
is commutative. If $M$ is an $A$-module, the diagram
\[
\begin{tikzcd}[column sep=huge,row sep=large]
S'^{-1}(M_{[\varphi]}) \arrow[r,"\sigma"] \arrow[d,"\rho^{T',S'}"']
  & (S^{-1}M)_{[\varphi^{S'}]} \arrow[d,"\rho^{T,S}"] \\
T'^{-1}(M_{[\varphi]}) \arrow[r,"\sigma"']
  & (T^{-1}M)_{[\varphi^{T'}]}
\end{tikzcd}
\]
is commutative. Finally, if $N'$ is an $A'$-module, the diagram
\[
\begin{tikzcd}[column sep=huge,row sep=large]
(S'^{-1}N')\otimes_{S'^{-1}A'}(S^{-1}A)_{[\varphi^{S'}]}
  \arrow[r,"\tau","{\sim}"'] \arrow[d]
  & S^{-1}(N'\otimes_{A'}A_{[\varphi]}) \arrow[d,"\rho^{T,S}"] \\
(T'^{-1}N')\otimes_{T'^{-1}A'}(T^{-1}A)_{[\varphi^{T'}]}
  \arrow[r,"\tau"',"{\sim}"]
  & T^{-1}(N'\otimes_{A'}A_{[\varphi]})
\end{tikzcd}
\]
is commutative, the left vertical arrow being obtained by applying
$\rho_{N'}^{T',S'}$ to $S'^{-1}N'$ and $\rho_A^{T,S}$ to $S^{-1}A$.
\end{env}

\oldpage[0\textsubscript{I}]{19}

\begin{env}[1.5.7]
\label{0.1.5.7}
Let $A''$ be a third ring, $\varphi':A''\to A'$ a ring
homomorphism, and $S''$ a multiplicative subset of $A''$ such that
$\varphi'(S'')\subset S'$. Set $\varphi''=\varphi\circ\varphi'$; then we have
\[
  \varphi''^{S''}=\varphi^{S'}\circ\varphi'^{S''}.
\]
Let $M$ be an $A$-module; obviously
$M_{[\varphi'']}=(M_{[\varphi]})_{[\varphi']}$; if $\sigma'$ and $\sigma''$ are
the homomorphisms defined from $\varphi'$ and $\varphi''$ as $\sigma$
is defined in (1.5.2) from $\varphi$, we have the transitivity formula
\[
  \sigma''=\sigma\circ\sigma'.
\]
Finally, let $N''$ be an $A''$-module; the $A$-module
$N''\otimes_{A''}A_{[\varphi'']}$ is canonically identified with
\[
  (N''\otimes_{A''}A'_{[\varphi']})\otimes_{A'}A_{[\varphi]},
\]
and similarly, the $S^{-1}A$-module
\[
  (S''^{-1}N'')\otimes_{S''^{-1}A''}
  (S^{-1}A)_{[\varphi''^{S''}]}
\]
is canonically identified with
\[
  \bigl((S''^{-1}N'')\otimes_{S''^{-1}A''}
  (S'^{-1}A')_{[\varphi'^{S''}]}\bigr)
  \otimes_{S'^{-1}A'}(S^{-1}A)_{[\varphi^{S'}]}.
\]
With these identifications, if $\tau'$ and $\tau''$ are the isomorphisms defined
from $\varphi'$ and $\varphi''$ as $\tau$ is defined in (1.5.4) from
$\varphi$, we have the transitivity formula
\[
  \tau''=\tau\circ(\tau'\otimes 1).
\]
\end{env}

\begin{env}[1.5.8]
\label{0.1.5.8}
Let $A$ be a subring of a ring $B$; for every \emph{minimal} prime
ideal $\mathfrak p$ of $A$, there exists a minimal prime ideal
$\mathfrak q$ of $B$ such that $\mathfrak p=A\cap\mathfrak q$. Indeed,
$A_{\mathfrak p}$ is a subring of $B_{\mathfrak p}$ (1.3.2) and has
only one prime ideal $\mathfrak p'$ (1.2.6); since $B_{\mathfrak p}$ is not
reduced to $0$, it has at least one prime ideal $\mathfrak q'$, and we
necessarily have $\mathfrak q'\cap A_{\mathfrak p}=\mathfrak p'$; the prime ideal
$\mathfrak q_1$ of $B$ which is the inverse image of $\mathfrak q'$ is therefore such that
$\mathfrak q_1\cap A=\mathfrak p$, and \emph{a fortiori} we have
$\mathfrak q\cap A=\mathfrak p$ for every minimal prime ideal $\mathfrak q$
of $B$ contained in $\mathfrak q_1$.
\end{env}

\subsection{Identification of the module $M_f$ with an inductive limit}
\label{subsection:0.1.6}

\begin{env}[1.6.1]
\label{0.1.6.1}
Let $M$ be an $A$-module and $f$ an element of $A$. Consider a sequence
$(M_n)$ of $A$-modules, all identical to $M$, and for every pair of integers
$m\leq n$, let $\varphi_{nm}$ be the homomorphism $z\mapsto f^{n-m}z$ from $M_m$
to $M_n$; it is immediate that $((M_n),(\varphi_{nm}))$ is an \emph{inductive
system} of $A$-modules; let $N=\varinjlim M_n$ be the inductive limit of this
system. We shall define a canonical \emph{functorial} $A$-isomorphism of
$N$ onto $M_f$. To this end, note that for every $n$,
$\theta_n:z\mapsto z/f^n$ is an $A$-homomorphism from $M=M_n$ to $M_f$, and
it follows from the definitions that $\theta_n\circ\varphi_{nm}=\theta_m$ for
$m\leq n$. There therefore exists an $A$-homomorphism $\theta:N\to M_f$ such that, if
$\varphi_n$ denotes the canonical homomorphism $M_n\to N$, we have
$\theta_n=\theta\circ\varphi_n$ for every $n$. Since by hypothesis every element
of $M_f$ is of the form $z/f^n$ for at least one $n$, it is clear that
$\theta$ is surjective. On the other hand, if $\theta(\varphi_n(z))=0$, in other words
$z/f^n=0$, there exists an integer $k>0$ such that $f^kz=0$, hence
$\varphi_{n+k,n}(z)=0$, which implies $\varphi_n(z)=0$. We can therefore
identify $M_f$ and $\varinjlim M_n$ by means of $\theta$.
\end{env}

\begin{env}[1.6.2]
\label{0.1.6.2}
Let us now write $M_{f,n}$, $\varphi_{nm}^f$ and $\varphi_n^f$ instead of
$M_n$, $\varphi_{nm}$ and $\varphi_n$. Let $g$ be a second element of $A$. Since
$f^n$ divides $f^ng^n$, we have a functorial homomorphism
\[
  \rho_{fg,f}:M_f\longrightarrow M_{fg}\qquad\text{(1.4.1 and 1.4.3)};
\]

\oldpage[0\textsubscript{I}]{20}

if $M_f$ and $M_{fg}$ are identified with $\varinjlim M_{f,n}$ and
$\varinjlim M_{fg,n}$ respectively, $\rho_{fg,f}$ is identified with the
\emph{inductive limit} of the maps
$\rho_{fg,f}^n:M_{f,n}\to M_{fg,n}$ defined by
$\rho_{fg,f}^n(z)=g^nz$. Indeed, this follows immediately from the
commutativity of the diagram
\[
\begin{tikzcd}[column sep=huge,row sep=large]
M_{f,n} \arrow[r,"\rho_{fg,f}^n"] \arrow[d,"\varphi_n^f"']
  & M_{fg,n} \arrow[d,"\varphi_n^{fg}"] \\
M_f \arrow[r,"\rho_{fg,f}"]
  & M_{fg}.
\end{tikzcd}
\]
\end{env}

\subsection{Support of a module}
\label{subsection:0.1.7}

\begin{env}[1.7.1]
\label{0.1.7.1}
Given an $A$-module $M$, the \emph{support} of $M$, denoted
$\operatorname{Supp}(M)$, is the set of prime ideals $\mathfrak p$ of $A$
such that $M_{\mathfrak p}\neq0$. For $M=0$, it is necessary and sufficient that
$\operatorname{Supp}(M)=\varnothing$, for if $M_{\mathfrak p}=0$ for every
$\mathfrak p$, the annihilator of an element $x\in M$ cannot be contained in
any prime ideal of $A$, and is therefore the whole of $A$.
\end{env}

\begin{env}[1.7.2]
\label{0.1.7.2}
If $0\to N\to M\to P\to0$ is an exact sequence of $A$-modules, we have
\[
  \operatorname{Supp}(M)=\operatorname{Supp}(N)\cup\operatorname{Supp}(P)
\]
since for every prime ideal $\mathfrak p$ of $A$ the sequence
$0\to N_{\mathfrak p}\to M_{\mathfrak p}\to P_{\mathfrak p}\to0$ is exact
(1.3.2), and for $M_{\mathfrak p}=0$ it is necessary and sufficient that
$N_{\mathfrak p}=P_{\mathfrak p}=0$.
\end{env}

\begin{env}[1.7.3]
\label{0.1.7.3}
If $M$ is the sum of a family $(M_\lambda)$ of submodules,
$M_{\mathfrak p}$ is the sum of the $(M_\lambda)_{\mathfrak p}$ for every prime
ideal $\mathfrak p$ of $A$ (1.3.3 and 1.3.2), hence
\[
  \operatorname{Supp}(M)=\bigcup_\lambda\operatorname{Supp}(M_\lambda).
\]
\end{env}

\begin{env}[1.7.4]
\label{0.1.7.4}
If $M$ is an $A$-module \emph{of finite type}, $\operatorname{Supp}(M)$ is
the set of prime ideals \emph{containing the annihilator of $M$}. Indeed,
if $M$ is monogenic and generated by $x$, to say that $M_{\mathfrak p}=0$ means
that there exists $s\notin\mathfrak p$ such that $s.x=0$, that is, that $\mathfrak p$ does not
contain the annihilator of $x$. If now $M$ has a finite system
$(x_i)_{1\leq i\leq n}$ of generators and $\mathfrak a_i$ is the annihilator
of $x_i$, it follows from (1.7.3) that $\operatorname{Supp}(M)$ is the set of
the $\mathfrak p$ containing one of the $\mathfrak a_i$, or, what amounts to the same,
the set of the $\mathfrak p$ containing $\mathfrak a=\bigcap_i\mathfrak a_i$,
which is the annihilator of $M$.
\end{env}

\begin{env}[1.7.5]
\label{0.1.7.5}
If $M$ and $N$ are two $A$-modules \emph{of finite type}, we have
\[
  \operatorname{Supp}(M\otimes_A N)
  =\operatorname{Supp}(M)\cap\operatorname{Supp}(N).
\]
We must show that if $\mathfrak p$ is a prime ideal of $A$, the
condition $M_{\mathfrak p}\otimes_{A_{\mathfrak p}}N_{\mathfrak p}\neq0$ is
equivalent to ``$M_{\mathfrak p}\neq0$ and $N_{\mathfrak p}\neq0$'' (taking
(1.3.4) into account). In other words, we must show that if $P$, $Q$ are two
modules of finite type over a \emph{local} ring $B$, not reduced to $0$, then
$P\otimes_B Q\neq0$. Let $\mathfrak m$ be the maximal ideal of $B$. By
Nakayama's lemma, the vector spaces $P/\mathfrak mP$ and
$Q/\mathfrak mQ$ are not reduced to $0$, hence neither is their
tensor product
\[
  (P/\mathfrak mP)\otimes_{B/\mathfrak m}(Q/\mathfrak mQ)
  =(P\otimes_B Q)\otimes_B(B/\mathfrak m),
\]
whence the conclusion.

In particular, if $M$ is an $A$-module of finite type and $\mathfrak a$ an ideal
of $A$, $\operatorname{Supp}(M/\mathfrak aM)$ is the set of prime
ideals containing both $\mathfrak a$ and the annihilator $\mathfrak n$ of
$M$ (1.7.4), in other words the set of prime ideals containing
$\mathfrak a+\mathfrak n$.
\end{env}

\oldpage[0\textsubscript{I}]{21}

